% ============================================================
%  Expression de Besoin / Problem Statement
%  Project: AEGIS Entropy — AI-Based Port Scan Detection
%  Updated: June 2026 (aligned with final report)
% ============================================================
\documentclass[12pt,a4paper]{article}

\usepackage[utf8]{inputenc}
\usepackage[T1]{fontenc}
\usepackage[english]{babel}
\usepackage{geometry}
\usepackage{enumitem}
\usepackage{booktabs}
\usepackage{array}
\usepackage{tabularx}
\usepackage{xcolor}
\usepackage{titlesec}
\usepackage{fancyhdr}
\usepackage{hyperref}

\geometry{margin=2.5cm}

% --- Colors ---
\definecolor{sectionblue}{RGB}{15, 23, 42}
\definecolor{accentcyan}{RGB}{6, 182, 212}

\titleformat{\section}
  {\large\bfseries\color{sectionblue}}
  {\thesection.}{0.5em}{}[\vspace{-0.5em}\textcolor{accentcyan}{\rule{\textwidth}{0.4pt}}]
\titleformat{\subsection}
  {\normalsize\bfseries\color{sectionblue}}
  {\thesubsection.}{0.5em}{}

\pagestyle{fancy}
\fancyhf{}
\fancyhead[L]{\small\textcolor{gray}{Expression de Besoin — AEGIS Entropy}}
\fancyhead[R]{\small\textcolor{gray}{AI Module — S2}}
\fancyfoot[C]{\small\thepage}
\renewcommand{\headrulewidth}{0.4pt}

\newcolumntype{L}[1]{>{\raggedright\arraybackslash}p{#1}}
\newcolumntype{C}[1]{>{\centering\arraybackslash}p{#1}}

\begin{document}

\begin{center}
  {\Large\bfseries Expression de Besoin / Probl\'{e}matique}\\[0.6em]
  {\large AEGIS Entropy}\\[0.3em]
  {\large D\'{e}tection de Port Scan et Reconnaissance R\'{e}seau\\[0.2em]
   par Apprentissage Automatique}\\[1.2em]
  \begin{tabular}{rl}
    \textbf{Module :}  & Intelligence Artificielle --- Semestre 2 \\
    \textbf{\'{E}quipe :} & 5 membres (Ing\'{e}nierie --- ENSA K\'{e}nitra) \\
    \textbf{Encadrant :}  & Pr.\ Aniss Moumen \\
    \textbf{Date :}    & Juin 2026 \\
  \end{tabular}
\end{center}

\vspace{1em}

% ============================================================
\section{Contexte et situation actuelle}

Dans le domaine de la s\'{e}curit\'{e} des r\'{e}seaux, la d\'{e}tection pr\'{e}coce des activit\'{e}s
de reconnaissance constitue un enjeu critique. Le \textit{Port Scanning} --- technique par
laquelle un attaquant sonde syst\'{e}matiquement les ports d'un r\'{e}seau pour identifier les
h\^{o}tes actifs, les services expos\'{e}s et les vuln\'{e}rabilit\'{e}s exploitables ---
repr\'{e}sente la premi\`{e}re phase de la quasi-totalit\'{e} des cyberattaques document\'{e}es
sous le mod\`{e}le \textit{Cyber Kill Chain}.

Actuellement, les infrastructures r\'{e}seau s'appuient principalement sur des pare-feux \`{a}
r\`{e}gles statiques et des syst\`{e}mes de d\'{e}tection d'intrusion (IDS) bas\'{e}s sur des
signatures, tels que Snort et Suricata.
Ces outils fonctionnent par reconnaissance de motifs connus et appliquent des seuils fixes.

Cette approche engendre plusieurs lacunes majeures :

\begin{itemize}[nosep,leftmargin=1.5em]
  \item \textbf{C\'{e}cit\'{e} face aux scans furtifs et lents :}
        les attaquants espacent volontairement leurs sondes pour rester sous les seuils
        de d\'{e}tection, rendant les r\`{e}gles statiques inefficaces.
  \item \textbf{Taux de faux positifs \'{e}lev\'{e} :}
        le trafic l\'{e}gitime --- mises \`{a} jour logicielles, d\'{e}couverte de services,
        outils de supervision --- d\'{e}clenche fr\'{e}quemment les m\^{e}mes signatures
        que les scans malveillants.
  \item \textbf{Absence de profilage comportemental :}
        les syst\`{e}mes actuels traitent chaque paquet isol\'{e}ment, sans analyser
        les motifs temporels et statistiques du trafic par adresse source sur des
        fen\^{e}tres temporelles glissantes.
  \item \textbf{Incapacit\'{e} \`{a} classifier le type de scan :}
        un IDS classique ne distingue pas un \textit{SYN Stealth Scan} d'un
        \textit{UDP Scan}, privant l'analyste d'information tactique essentielle.
\end{itemize}

Ces limitations impactent directement les \'{e}quipes SOC en termes de temps de r\'{e}action,
de fiabilit\'{e} des alertes et de capacit\'{e} \`{a} contenir une attaque avant qu'elle
ne progresse vers des phases plus destructrices.

% ============================================================
\section{Probl\'{e}matique}

Face \`{a} ces constats, il devient n\'{e}cessaire de repenser le m\'{e}canisme de d\'{e}tection
afin de \textbf{d\'{e}passer les limites des approches statiques} et doter le r\'{e}seau
d'une capacit\'{e} d'analyse intelligente, adaptative et temps-r\'{e}el.

L'objectif principal est de concevoir un syst\`{e}me capable d'apprendre automatiquement
les motifs comportementaux du trafic r\'{e}seau --- normal comme malveillant --- afin de
d\'{e}tecter de mani\`{e}re autonome les activit\'{e}s de reconnaissance, y compris
les variantes furtives.

Dans ce contexte, plusieurs questions se posent :

\begin{enumerate}[nosep,leftmargin=1.5em]
  \item Quelles caract\'{e}ristiques (\textit{features}) du trafic r\'{e}seau discriminent
        efficacement un port scan d'un comportement l\'{e}gitime ?
  \item Comment un mod\`{e}le d'apprentissage automatique peut-il d\'{e}tecter des
        scans distribu\'{e}s dans le temps, l\`{a} o\`{u} les seuils statiques \'{e}chouent ?
  \item Quelles donn\'{e}es \'{e}tiquet\'{e}es et repr\'{e}sentatives faut-il fournir
        au mod\`{e}le pour garantir une classification fiable avec un taux minimal
        de faux positifs ?
\end{enumerate}

\vspace{0.5em}
\noindent La probl\'{e}matique se formule ainsi :

\begin{quote}
\itshape
\textit{Comment peut-on d\'{e}tecter et classifier les attaques de type PortScan
\`{a} partir de donn\'{e}es de trafic r\'{e}seau (pcap/csv) afin de prot\'{e}ger
l'infrastructure et alerter les \'{e}quipes SOC en temps r\'{e}el ?}
\end{quote}

Cette probl\'{e}matique a \'{e}t\'{e} valid\'{e}e par quatre entretiens semi-directifs
conduits aupr\`{e}s d'\'{e}tudiants en ing\'{e}nierie r\'{e}seaux et cybers\'{e}curit\'{e}.
Les personnes interrog\'{e}es ont confirm\'{e} que la surveillance manuelle est insoutenable,
que les faux positifs repr\'{e}sentent un point de douleur majeur, et qu'un syst\`{e}me
capable d'apprentissage comportemental avec r\'{e}ponse autonome gradu\'{e}e r\'{e}pond
\`{a} un besoin r\'{e}el.

% ============================================================
\section{Solution propos\'{e}e}

Pour r\'{e}soudre cette probl\'{e}matique, nous adoptons une approche fond\'{e}e sur
\textbf{l'intelligence artificielle}, et plus sp\'{e}cifiquement sur
\textbf{l'apprentissage automatique supervis\'{e} et non supervis\'{e}}.

Le syst\`{e}me, baptis\'{e} \textbf{AEGIS Entropy} (\textit{Adaptive Entropy-based
Gateway for Intrusion Suppression}), combine deux orientations :

\begin{enumerate}
  \item \textbf{Orientation aide \`{a} la d\'{e}cision :}
        un tableau de bord SOC temps-r\'{e}el qui affiche les scores de confiance
        de classification, les scanners actifs, les contacts honeypot et les IP bloqu\'{e}es,
        permettant aux analystes de comprendre et de valider les d\'{e}cisions du syst\`{e}me.
  \item \textbf{Orientation automatisation :}
        un moteur de d\'{e}fense active qui redirige automatiquement les attaquants vers des
        leurres, mute la surface d'attaque expos\'{e}e (MTD), et bloque les adresses sources
        malveillantes sans intervention humaine.
\end{enumerate}

\subsection{Entr\'{e}es et sorties}

\begin{tabular}{@{}l@{\hspace{1em}}p{11cm}@{}}
  \textbf{Entr\'{e}es :}
    & Vecteurs de 11 caract\'{e}ristiques extraits du trafic r\'{e}seau :
      Destination Port, Flow Duration, Total Fwd Packets, SYN/RST/ACK Flag Count,
      Flow IAT Mean, Bwd Packet Length Mean, Init\_Win\_bytes\_forward,
      et deux placeholders (\texttt{shadow\_node\_interaction},
      \texttt{mtd\_port\_delta}) activ\'{e}s en production. \\[0.5em]
  \textbf{Sorties :}
    & Classification binaire \textsc{Benign} / \textsc{PortScan} avec score de
      confiance ; une extension multi-classe (SYN / UDP / ACK / XMAS / Connect)
      est pr\'{e}vue en perspective. \\[0.5em]
  \textbf{Mod\`{e}les :}
    & Random Forest (classifieur principal, 100 arbres, bagging)
      + XGBoost (benchmark, gradient boosting, 100 estimateurs)
      + Isolation Forest (d\'{e}tection non supervis\'{e}e des scans lents).
\end{tabular}

\subsection{M\'{e}thodologie}

Le projet suit le cadre \textbf{CRISP-DM} (\textit{Cross-Industry Standard Process for
Data Mining}) : compr\'{e}hension du besoin m\'{e}tier, compr\'{e}hension des donn\'{e}es,
pr\'{e}paration des donn\'{e}es, mod\'{e}lisation, \'{e}valuation, et d\'{e}ploiement.
Des contr\^{o}les stricts de non-fuite de donn\'{e}es (\textit{anti-leakage}) sont
appliqu\'{e}s : le StandardScaler est ajust\'{e} exclusivement sur l'ensemble
d'entra\^{i}nement, et SMOTE est appliqu\'{e} uniquement apr\`{e}s la s\'{e}paration
train/test.

\subsection{Donn\'{e}es}

Le jeu de donn\'{e}es principal est \textbf{CICIDS2017}, sous-ensemble
\textit{Friday-WorkingHours-Afternoon-PortScan}, contenant 286\,467 flux r\'{e}seau
(158\,930 PortScan, 127\,537 BENIGN) avec 79 colonnes brutes, r\'{e}duites \`{a}
11 caract\'{e}ristiques apr\`{e}s s\'{e}lection.

% ============================================================
\section{Synth\`{e}se des r\'{e}sultats}

Le syst\`{e}me a \'{e}t\'{e} enti\`{e}rement impl\'{e}ment\'{e} et valid\'{e} sur
les 57\,294 flux de l'ensemble de test.

\begin{table}[h!]
\centering
\caption{Performances obtenues sur l'ensemble de test CICIDS2017}
\begin{tabular}{lcccc}
\toprule
\textbf{Mod\`{e}le} & \textbf{Accuracy} & \textbf{F1-Score} & \textbf{FPR} & \textbf{AUC-ROC} \\
\midrule
Random Forest & 99.91\% & 99.92\% & 0.11\% & 0.9997 \\
XGBoost       & 99.91\% & 99.92\% & 0.11\% & 0.9999 \\
Isolation Forest & 24.63\% & 9.46\% & 53.53\% & --- \\
\bottomrule
\end{tabular}
\end{table}

\noindent Les deux mod\`{e}les supervis\'{e}s d\'{e}passent largement les cibles initiales
(F1 $\geq$ 88\%, Accuracy $\geq$ 90\%, FPR $<$ 6\%). L'Isolation Forest \'{e}choue
sur ce jeu de donn\'{e}es en raison du \textbf{paradoxe de l'anomalie majoritaire} :
PortScan constitue 55.48\% des donn\'{e}es --- la classe majoritaire --- ce qui viole
l'hypoth\`{e}se de raret\'{e} des anomalies. En production, l'Isolation Forest sera
r\'{e}-entra\^{i}n\'{e} en mode semi-supervis\'{e} sur le trafic BENIGN uniquement.

Le d\'{e}tail complet du pipeline de donn\'{e}es, des mod\`{e}les, de l'architecture
de d\'{e}fense et de l'int\'{e}gration est document\'{e} dans le rapport final
(\texttt{AEGIS.pdf}, 9 chapitres + annexes).

\end{document}
